%
% vox-notify.tex
%
% Z Specification for vox notification state machine.
% Formal model of the notify/speak configuration
% and the commands that transition between states.
%

\documentclass[a4paper,10pt,fleqn]{article}
\usepackage[margin=1in]{geometry}
\usepackage{fuzz}
\usepackage[colorlinks=true,linkcolor=blue,citecolor=blue,urlcolor=blue]{hyperref}
\newcommand{\ok}{$\surd$}

\begin{document}

\title{Vox Notification State Machine: A Z Specification}
\author{Formal Model of notify/speak Configuration}
\date{March 2026}
\hypersetup{
  pdftitle={Vox Notification State Machine: A Z Specification},
  pdfauthor={Punt Labs},
  pdfsubject={Z Specification},
  pdfcreator={fuzz/probcli}
}
\maketitle

\tableofcontents
\newpage

% ===================================================================
\section{Introduction}
% ===================================================================

Vox has two independent configuration axes that control audio
feedback during a Claude Code session:

\begin{itemize}
  \item \textbf{notify} --- which events trigger feedback:
    off ($nOff$), on-demand ($nOn$, stop and prompts only),
    or continuous ($nCont$, real-time signals too).
  \item \textbf{speak} --- what form feedback takes:
    chime ($sChime$) or voice ($sVoice$, synthesized speech).
\end{itemize}

The specification models the nine-cell state space (seven reachable), the commands
that transition between states, and the constraints imposed by
the current implementation. Its purpose is to identify
unreachable states and missing commands.

% ===================================================================
\section{Basic Types}
% ===================================================================

\begin{zed}
  [VOICE]
\end{zed}

$VOICE$ represents a voice identifier (e.g.\ ``matilda'', ``roger'').
It has no internal structure in this model.

% ===================================================================
\section{Free Types}
% ===================================================================

\begin{zed}
  NotifyMode ::= nOff | nOn | nCont
\end{zed}

\begin{zed}
  SpeakMode ::= sUnset | sChime | sVoice
\end{zed}

$sUnset$ represents the state before the user has made an
explicit choice. Commands that enable notifications initialize
$speak$ to $sVoice$ when it is $sUnset$, but preserve the
user's choice once set.

% ===================================================================
\section{State}
% ===================================================================

\begin{schema}{VoxState}
  notify : NotifyMode \\
  speak : SpeakMode \\
  voice : \power VOICE
  \where
  \# voice \leq 1
\end{schema}

The $voice$ field is modeled as a set of at most one element ---
an optional value. When empty, the provider default is used.

% ===================================================================
\section{Initialization}
% ===================================================================

\begin{schema}{Init}
  VoxState'
  \where
  notify' = nOff \\
  speak' = sUnset \\
  voice' = \emptyset
\end{schema}

The initial state has notifications off and speak uninitialized.
The first command that enables notifications will initialize
$speak$ to $sVoice$.

% ===================================================================
\section{Observations}
% ===================================================================

These schemas model the three feedback consumers. They do not
change state --- they observe it.

\subsection{Stop/prompt feedback}

Fires when $notify \neq nOff$. Output depends on $speak$.
When $speak = sUnset$, the system behaves as $sVoice$
(voice is the default).

\begin{schema}{StopFires}
  \Xi VoxState
  \where
  notify \neq nOff
\end{schema}

\begin{schema}{StopChime}
  \Xi VoxState
  \where
  notify \neq nOff \\
  speak = sChime
\end{schema}

\begin{schema}{StopVoice}
  \Xi VoxState
  \where
  notify \neq nOff \\
  speak \neq sChime
\end{schema}

\subsection{Watcher (real-time signals)}

Fires only when $notify = nCont$. Output depends on $speak$.

\begin{schema}{WatcherFires}
  \Xi VoxState
  \where
  notify = nCont
\end{schema}

\begin{schema}{WatcherChime}
  \Xi VoxState
  \where
  notify = nCont \\
  speak = sChime
\end{schema}

\begin{schema}{WatcherVoice}
  \Xi VoxState
  \where
  notify = nCont \\
  speak \neq sChime
\end{schema}

% ===================================================================
\section{Commands}
% ===================================================================

These are the state transitions triggered by user commands.

\begin{center}
\begin{tabular}{l|l|l}
\textbf{User command} & \textbf{Operation} & \textbf{Effect} \\
\hline
\texttt{/vox y}          & $VoxOn$   & $notify \to nOn$; $speak \to sVoice$ if $sUnset$ \\
\texttt{/vox n}          & $VoxOff$  & $notify \to nOff$; $speak$ unchanged \\
\texttt{/vox c}          & $VoxCont$ & $notify \to nCont$; $speak \to sVoice$ if $sUnset$ \\
\texttt{/unmute}         & $Unmute$  & $speak \to sVoice$; $notify$ unchanged; optional voice \\
\texttt{/unmute @name}   & $Unmute$  & Same as \texttt{/unmute} with $v? = \{name\}$ \\
\texttt{/mute}           & $Mute$    & $speak \to sChime$; $notify$ unchanged \\
\end{tabular}
\end{center}

The two axes are controlled by separate commands.
\texttt{/vox y\textbar n\textbar c} controls the $notify$ axis
and initializes $speak \to sVoice$ on first use (when $sUnset$),
but preserves the user's explicit choice thereafter.
\texttt{/unmute} and \texttt{/mute} control the $speak$ axis
independently. Chime-only states require an explicit
\texttt{/mute}. The design invariant: \textbf{first activation
implies voice; subsequent activations respect the user's choice}.

The observation schemas are triggered by system events, not
user commands:

\begin{center}
\begin{tabular}{l|l|l}
\textbf{System event} & \textbf{Observation} & \textbf{When it fires} \\
\hline
Task completion      & $StopChime$ / $StopVoice$       & $notify \neq nOff$ \\
Permission prompt    & $StopChime$ / $StopVoice$       & $notify \neq nOff$ \\
Real-time signal     & $WatcherChime$ / $WatcherVoice$ & $notify = nCont$ only \\
\end{tabular}
\end{center}

\subsection{VoxOn ($\mathord{/vox~y}$)}

Sets $notify = nOn$. Initializes $speak$ to $sVoice$ if
unset; preserves the user's choice otherwise.

\begin{schema}{VoxOn}
  \Delta VoxState
  \where
  notify' = nOn \\
  speak' = \IF speak = sUnset \THEN sVoice \ELSE speak \\
  voice' = voice
\end{schema}

\subsection{VoxOff ($\mathord{/vox~n}$)}

Disables all notifications.

\begin{schema}{VoxOff}
  \Delta VoxState
  \where
  notify' = nOff \\
  speak' = speak \\
  voice' = voice
\end{schema}

\subsection{VoxCont ($\mathord{/vox~c}$)}

Enables continuous mode. Initializes $speak$ to $sVoice$
if unset; preserves the user's choice otherwise.

\begin{schema}{VoxCont}
  \Delta VoxState
  \where
  notify' = nCont \\
  speak' = \IF speak = sUnset \THEN sVoice \ELSE speak \\
  voice' = voice
\end{schema}

\subsection{Unmute ($\mathord{/unmute}$)}

Sets $speak = sVoice$ without changing $notify$.
Optionally sets a session voice. Symmetric inverse
of $Mute$.

\begin{schema}{Unmute}
  \Delta VoxState \\
  v? : \power VOICE
  \where
  \# v? \leq 1 \\
  notify' = notify \\
  speak' = sVoice \\
  voice' = \IF v? = \emptyset \THEN voice \ELSE v?
\end{schema}

\subsection{Mute ($\mathord{/mute}$)}

Switches to chime mode. Does not change $notify$.

\begin{schema}{Mute}
  \Delta VoxState
  \where
  notify' = notify \\
  speak' = sChime \\
  voice' = voice
\end{schema}

% ===================================================================
\section{Reachability Analysis}
% ===================================================================

The state space has nine cells (3 notify $\times$ 3 speak),
but $sUnset$ is a transient initial state --- any command that
touches the $speak$ axis or enables notifications moves out of
it permanently.

\begin{center}
  \begin{tabular}{l|c|c|c}
    & $sUnset$ & $sChime$ & $sVoice$ \\
    \hline
    $nOff$  & cell 0 & cell 1 & cell 2 \\
    $nOn$   & --- & cell 3 & cell 4 \\
    $nCont$ & --- & cell 5 & cell 6 \\
  \end{tabular}
\end{center}

Cells marked --- are unreachable: enabling notifications ($nOn$
or $nCont$) initializes $sUnset \to sVoice$, so $speak$ is never
$sUnset$ when notifications are active.

\subsection{Full transition table}

Every cell lists the target cell for each command.

\begin{center}
  \begin{tabular}{l|c|c|c|c|c}
    \textbf{From} & \texttt{/vox y} & \texttt{/vox n} & \texttt{/vox c} & \texttt{/unmute} & \texttt{/mute} \\
    \hline
    0 ($nOff/sUnset$)  & 4 & 0 & 6 & 2 & 1 \\
    1 ($nOff/sChime$)  & 3 & 1 & 5 & 2 & 1 \\
    2 ($nOff/sVoice$)  & 4 & 2 & 6 & 2 & 1 \\
    3 ($nOn/sChime$)   & 3 & 1 & 5 & 4 & 3 \\
    4 ($nOn/sVoice$)   & 4 & 2 & 6 & 4 & 3 \\
    5 ($nCont/sChime$) & 3 & 1 & 5 & 6 & 5 \\
    6 ($nCont/sVoice$) & 4 & 2 & 6 & 6 & 5 \\
  \end{tabular}
\end{center}

Note: from cells 1, 3, 5 (chime), \texttt{/vox y} and \texttt{/vox c}
preserve $sChime$ because $speak \neq sUnset$. This is the key
design property: once the user has muted, enabling notifications
does not re-unmute them.

\subsection{Reachability}

Starting from $Init$ ($nOff$, $sUnset$ = cell 0):

\begin{itemize}
  \item \textbf{Cell 4} ($nOn/sVoice$): $VoxOn$ from cell 0. \ok{}
  \item \textbf{Cell 6} ($nCont/sVoice$): $VoxCont$ from cell 0. \ok{}
  \item \textbf{Cell 2} ($nOff/sVoice$): $Unmute$ from cell 0. \ok{}
  \item \textbf{Cell 1} ($nOff/sChime$): $Mute$ from cell 0. \ok{}
  \item \textbf{Cell 3} ($nOn/sChime$): $Mute$ then $VoxOn$. \ok{}
  \item \textbf{Cell 5} ($nCont/sChime$): $Mute$ then $VoxCont$. \ok{}
\end{itemize}

All seven reachable cells are accessible from $Init$ in at most
two commands. The $sUnset$ state is left on the first command
that touches either axis and is never re-entered.

The design has two invariants:
\begin{enumerate}
  \item \textbf{First activation implies voice}: $sUnset$ resolves
    to $sVoice$ on first enable.
  \item \textbf{User choice is sticky}: once $speak$ is $sChime$ or
    $sVoice$, only \texttt{/mute} and \texttt{/unmute} change it.
    Enabling notifications preserves the user's explicit choice.
\end{enumerate}

No missing commands.

\end{document}
